\FourSubsection{Inputs Common to All Periodic Boundaries}
\label{SEC:BC:Numerical:Periodic:Inputs:Common}

%This section discusses the inputs for specifying a pair of coupled periodic
%boundaries that are common to all periodic boundaries whether the periodicity
%is translational or rotational.
%
This section discusses inputs that can be used with both translational or
rotational periodic boundaries.
%
The \inptext{Periodic} subcomponent of \BCinpBC[nmlBC::BCInput]{} is used
to provide the information that will allow GFR to group two boundary conditions
into a matching pair of periodic boundaries, which in this example will be
\BCinptext[i]{} and \BCinptext[j]{}.
%
Both sides of a periodic pair must provide their respective information in the
input file.
%
%As you will see, the information required by opposing sides is the reverse\slash
%negative of each other, and could therefore probably be streamlined to reduce
%redundancies.
%%
%However, there was intentionality in how this was set up to help aid in error
%checking and making sure the user understands what they are doing in setting up
%periodic boundary conditions.
%
As you will see, the information required by opposing sides is the reverse\slash
negative of each other, and is used to aid error checking of the periodic
boundary setup.

%%%%%%%%%%%%%%%%%%%%%%%%%%%%%%%%%%%%%%%%%%%%%%%%%%%%%%%%%%%%%%%%%%%%%%%%%%%%%%%%

The two inputs used to define a matching periodic pair are
%
\begin{InputExample}
  \item[{\BCinptext[i]{Periodic\%Matching\_BC\_Name}}]
  \item[{\BCinptext[i]{Periodic\%Matching\_BC\_Idx}}]
\end{InputExample}
%
The values for these two inputs need to provide the respective information from
the opposing matching periodic boundary condition.
%
Using boundary conditions $i$ and $j$ defined at the beginning of
\sref{SEC:BC:Numerical:Periodic}, the inputs for boundary condition $i$ would
look like
%
\begin{InputExample}
  \item[{\BCinptext[i]{Name = "Face Side A"}}]
  \item[{\BCinptext[i]{Set\_BC\_Value = 81}}]
  \item[{\BCinptext[i]{Periodic\%Matching\_BC\_Name = "Periodic-2"}}]
  \item[{\BCinptext[i]{Periodic\%Matching\_BC\_Idx = }j}]
\end{InputExample}
%
and the inputs for boundary condition $j$ would look like
%
\begin{InputExample}
  \item[{\BCinptext[j]{Name = "Periodic-2"}}]
  \item[{\BCinptext[j]{Set\_BC\_Value = 81}}]
  \item[{\BCinptext[j]{Periodic\%Matching\_BC\_Name = "Face Side A"}}]
  \item[{\BCinptext[j]{Periodic\%Matching\_BC\_Idx = }i}]
\end{InputExample}
%
The remaining parts of \BCinptext{Periodic} are used to define the physical
orientation between the pair of periodic boundaries, and will be covered in the
next two sections.

%###############################################################################

